\subsection{The three dihedral genus-three covers: the complete fixed subgroup}

\begin{statement}{A.37}{Proved}{All first dihedral genus-three covers have exactly cusp fixed subgroups}
For each of the three curves
\[
 C=C_{4,a,b}:v^2=w^8+aw^4+b,
 \qquad (a,b)\in\{(0,-1),(6,1),(-6,1)\},
\]
let \(S\) be the four cusps of A.4. For every prime \(\ell\),
every packet section induced as in Part A, and every open subgroup
\(\Delta_0\) preserving the selected cover and cusp lift \(P\),
\[
 \left(\pi_1((C\setminus S)_{\bar K})^{(\ell)}\right)^{s_b(\Delta_0)}
       =I_{P,\ell}.                                      \tag{A.37.1}
\]
At odd primes the proof works for every cusp packet over a finite
number field of definition, with that field's own derived Galois
subgroup. No nonzero Kummer condition on the packet is imposed.

The pointed stable reduction at two has exactly two rational
components carrying the cusps. Each carries two cusps and meets
the rest of the fibre at a single node. The other components can
have different configurations for the three curves. The assertion
uses only these two marked leaves and does not identify the whole
pro-\(\ell\) group with its image under a covering map.
\end{statement}

\begin{proof}
\paragraph{A harmonic cyclic-four map.}
Choose \(s^2=b\) and put \(T=(v-s)/w^4\). Direct substitution gives
\[
 w^4=\frac{a-2sT}{T^2-1},\qquad v=s+Tw^4.               \tag{A.37.2}
\]
Thus \(\mu(w,v)=(iw,v)\) generates a cyclic-four quotient
to \(\mathbf P^1_T\), totally ramified at precisely the four
cusps, whose images are
\[
                 a/(2s),\quad\infty,\quad 1,\quad-1.
                                                               \tag{A.37.3}
\]
Every nonidentity power of \(\mu\), including
\(\sigma=\mu^2\), fixes precisely these four points:
on the affine chart its fixed points must have \(w=0\),
and at infinity the two branches \(v/w^4=\pm1\)
are individually fixed. There are no additional generic
ramification points to be marked.
For \(a=0\) this is the harmonic set \(\{0,\infty,1,-1\}\).
For \((a,s)=(6,1)\), the substitution \(z=(T-1)/2\)
identifies it with the same set; for \((a,s)=(-6,1)\), use
\(z=(T+1)/2\). In every case its cross-ratios are
\(-1,2,1/2\), so this four-pointed projective line has bad
marked reduction at two, after every finite extension.
Indeed four distinct reductions, after normalizing three
of them to \(0,1,\infty\), would make both the fourth
coordinate and its difference from one units. None of
these three cross-ratios has that property at two.

The original \(D_8\) deck group from A.4 is generated by
\[
 \alpha(w,v)=(iw,-v),\qquad
 j(w,v)=(c/w,-c^2v/w^4),\qquad c^4=b.
\]
It acts transitively on the four cusps and commutes with
\(\sigma=\alpha^2=\mu^2\). The quotient by \(\mu\) has
genus zero, while the quotient by \(\sigma\) has genus one.
Consequently their invariant subspaces on proper rational
\(\ell\)-adic homology have dimensions zero and two,
respectively. These dimensions persist in a compact-type
stable fibre.

\paragraph{The possible unpointed special fibres.}
All proper factors of the Jacobians in A.4 have integral
\(j\)-invariants reducing to zero at two: the relevant values
are \(1728,8000,10976,128\). After a finite extension they
therefore have good supersingular reduction. The Jacobian of
\(C\) has good reduction with \(2\)-rank zero. Its stable
fibre is of compact type, and each positive-genus component
has \(2\)-rank zero. An unmarked stable genus-three tree
can consequently have one genus-three component, one
genus-two and one elliptic component, or three elliptic
components. In the last case its only shapes are a chain
of three elliptic components or three elliptic tails
attached to a rational centre. There are no additional
rational vertices: stability gives each such vertex
valence at least three, and the tree has at most three
positive-genus vertices.

We need a short characteristic-two restriction. Every
nontrivial two-power automorphism of a smooth curve of
\(2\)-rank zero has exactly one geometric fixed point;
all nontrivial elements of a two-group have the same
fixed point. This is the Deuring--Shafarevich argument in
\cite[Lemma 3.2 and Remark 3.1]{fghGiuliettiKorchmaros}.
For a faithful cyclic-four action there is therefore
one ramification point. If \(u_1,u_2\) are its upper
jumps, then
\[
 u_1>0\text{ is odd},\qquad
 u_2=2u_1\quad\text{or}\quad u_2>2u_1\text{ is odd},
                                                               \tag{A.37.4}
\]
by \cite[Theorem 1.1(c)]{fghObusPries}.
The lower jumps are \(u_1,2u_2-u_1\), and the different
is \(u_1+2u_2+3\). Riemann--Hurwitz gives
\[
      g=4g_{\rm quot}+\frac{u_1+2u_2-3}{2}.             \tag{A.37.5}
\]
In particular genus three is impossible: a positive quotient
genus gives \(g\ge5\), while genus zero would require
\(u_1+2u_2=9\), contrary to (A.37.4).
For genus two the only possibility is
\(g_{\rm quot}=0,(u_1,u_2)=(1,3)\); the quotient by
the unique involution also has genus zero, since its
different is six.

In the smooth genus-three case the specialization of
\(\sigma\) is nonidentity: its homology fixed space
has dimension two, rather than six. Thus \(\mu\)
specializes to a faithful order-four automorphism, and
(A.37.5) rules this case out.

If the fibre has genera two and one, \(\mu\) preserves
both components. It must have order four on the genus-two
component: otherwise \(\sigma\) would act identically
there and have at least four homology invariants.
The preceding genus-two calculation shows that
\(\sigma\) has no invariants on that component.
It is therefore the identity on the elliptic component,
where \(\mu\) is nontrivial of order two. Each component's
unique \(\mu\)-fixed point is their joining node.
All four \(\mu\)-fixed cusp sections specialize to that node.

Suppose the fibre is an elliptic chain. If \(\mu\)
preserved each component, its action on the middle
component would be nontrivial, because \(\mu\) has no
homology invariants. It would fix two distinct attaching
points, contradicting the fixed-point restriction.
Hence \(\mu\) exchanges the two outer components.
On those two components \(\mu^2\) is nontrivial;
otherwise their exchanged homology spaces would supply
\(\mu\)-invariants. All \(\sigma\)-invariants therefore
come from the middle elliptic component, on which
\(\mu\) is an involution. Its unique fixed point is
smooth, since its two attaching points are exchanged.
All four cusps specialize to this one smooth point.

It remains to consider the three-tail star. If \(\mu\)
exchanges two tails, it acts on the rational centre as
a nonidentity involution. Its unique fixed point there
is the attaching node of the third tail. On that third
tail \(\mu\) is a nontrivial involution, again with
unique fixed point that node. All four cusps specialize
to the same node. If instead \(\mu\) preserves all
tails, it is the identity on the centre, since it fixes
all three attaching points. It acts with order two
on exactly one tail \(E_0\), and order four on the
other two: this is forced by the two-dimensional
\(\sigma\)-invariant space. Its only fixed point on
each tail is the attaching node. All cusp specializations
therefore lie on the rational centre, including possibly
at its nodes. The \(D_8\) action preserves \(E_0\), as
it commutes with \(\sigma\). Its action on the centre
embeds in the permutations of the other two attaching
points, so has order at most two. Transitivity on the
four cusps shows that their specialization set consists
of either one point or two points, with clusters of
four or two cusps respectively.

\paragraph{Separating the marked clusters.}
In every case above there are one or two cusp clusters,
each of even size, and each can be separated in a
rational branch containing no positive-genus component.
For a cluster at a smooth point this is the ordinary
marked resolution. For a cluster at a node, its
stabilizer in \(D_8\) preserves both adjacent components
and hence fixes the two ends of the metric edge.
It acts identically on the edge. Since it is transitive
on the cluster, all its cusp sections have the same
interior retraction point. After subdividing there,
the stabilizer fixes both attachments of the inserted
rational bridge. A nontrivial two-power automorphism
of a projective line in characteristic two has only
one geometric fixed point, so this bridge action is
trivial. The cluster thus meets a common smooth point
on that bridge, and its marked separation attaches a
rational branch containing no genus.

The same rooted-tree stability argument as in A.36
now applies within each cluster. Transitivity forces
all marked vertices to be leaves, each with at least
two marks. A two-point cluster gives a single two-mark
leaf. A four-point cluster gives either one four-mark
leaf or two two-mark leaves.

We exclude the four-mark leaf using the actual cyclic
quotient (A.37.2). The four cusps are precisely all its
ramification points. The stable marked model is therefore
\(\langle\mu\rangle\)-semistable in the sense of
\cite[Definition 3.1(c)]{fghBouwWewers}.
Its quotient is a semistable model of
\(\mathbf P^1_T\), and the four branch sections specialize
to distinct smooth points, by
\cite[Lemma 3.2(i)--(ii)]{fghBouwWewers}.
If a source component carried all four cusps, \(\mu\)
would fix these four points and hence act identically
on that component. Its image under the finite quotient
would be one component containing all four distinct
smooth branch marks. Marked stabilization preserves
that component and contracts every other component:
the target has genus zero and no marks elsewhere.
It would give good reduction of the four-pointed
projective line, contradicting (A.37.3).
This argument uses the semistable quotient theorem
with its marked-ramification hypotheses; it does not
identify the wild scheme quotient with the geometric
special-fibre quotient without justification.

There are consequently precisely two two-mark rational
leaves and no cusps on the remaining components.

\paragraph{The genuine first derivation and its kernel.}
Fix the selected leaf with cusps \(P,Q\), and denote
the other leaf's cusps by \(R_1,R_2\). Let
\(p,q,r\) be the classes of \(P,Q,R_1\) in ordinary
homology, put \(s_0=p+q\), and choose six handle
generators \(v_1,\ldots,v_6\). The last cusp class
is \(-s_0-r\). For a base point on the selected leaf,
the based edge-twist formula gives
\[
 \begin{split}
 D_0p&=D_0q=0,\\
 D_0v_j&=-a_j[s_0,v_j],\qquad a_j>0,\\
 D_0r&=-L[s_0,r],\qquad L>0.
 \end{split}
                                                        \tag{A.37.6}
\]
Here \(L\) is the length of the path between the two
marked leaves. For a handle on a positive-genus
component, \(a_j\) is the distance from the selected
leaf to the attachment of that component's path to
this leaf-to-leaf path. Each \(a_j\) is positive
because the selected component is a genus-zero leaf.
Every edge on this path has first boundary class
\(-s_0\), while edges off the path enclose only
genus and have zero first boundary class. This proves
(A.37.6) directly from
\cite[Definition 3.1.5 and Theorem 3.1.6]{fghBettsDogra}.

The ordinary-point to tangential-point comparison in
A.36 takes place inside the selected cusp's punctured
formal disc. It changes this genuine first derivation
only by \(\lambda\operatorname{ad}p\), for some
\(\lambda\in\mathbf Q_\ell\). This includes every
completed tangent packet. Choose a finite global field
\(L_0\) with \(\Gamma=\overline{[G_{L_0},G_{L_0}]}\subset\Delta_0\).
An actual commutator of local Frobenius at two with
tame inertia has first derivation a nonzero multiple
of
\[
                 D_\lambda=D_0+\lambda\operatorname{ad}p.
                                                               \tag{A.37.7}
\]
As in A.36, no value of \(\lambda\) is assumed.

The all-degree first-slot proof in A.36 applies
without change: on the free subalgebra generated by
\(p,s_0\), the derivation is \(\lambda\operatorname{ad}p\);
on every outside letter \(z\), it is
\([\lambda p+c_zs_0,z]\), with \(c_z\ne0\).
The first-slot operator \(\lambda L_p+c_zR_{s_0}\)
is injective. Hence, inside the ordinary free Lie algebra,
\[
                  \ker D_\lambda\subset\mathbb L(p,q).
                                                               \tag{A.37.8}
\]
This is an unbounded statement, including \(\lambda=0\).

\paragraph{The harmonic target detects the remaining leading algebra.}
After a projective change of \(T\), the four branch
points are \(\{0,\infty,1,-1\}\). The two maps
\[
                   z\longmapsto z^2,
          \qquad z\longmapsto\left(\frac{z-1}{z+1}\right)^2
                                                               \tag{A.37.9}
\]
are degree-two maps to the tripod, each having exactly
these four preimages of \(\{0,1,\infty\}\).
The first folds \(\{1,-1\}\); the second folds
\(\{0,\infty\}\). Thus at least one sends \(P,Q\)
to distinct target cusps. Compose it with (A.37.2).
The result is an actual degree-eight etale morphism
of punctured curves, with selected cusp indices
\(e_P,e_Q\in\{4,8\}\).
On first homology it sends \(p,q\) to nonzero
multiples of two distinct tripod cusp classes.
It is therefore injective on \(\mathbb L(p,q)\)
in every degree over \(\mathbf Q_\ell\).

If \(g\) is an actual \(\Delta_0\)-fixed element,
A.34 puts its image in the target's selected cusp
inertia. At odd \(\ell\), \(e_P\) is a unit in
\(\mathbf Z_\ell\), so subtract an actual integral
power of the selected cusp generator to obtain an
element \(h\) which is still fixed and maps to one.
If \(h\ne1\), its first nonzero ordinary lower-central
class survives rationally, is killed by (A.37.7),
and hence lies in (A.37.8). Its image cannot vanish
under the target map, a contradiction. Thus \(h=1\).
The opposite inclusion follows from trivial cyclotomic
action on \(\Delta_0\). Finally, \(\ell=2\) is A.22,
since each original \(D_8\) cover has two-group
Galois closure. This proves (A.37.1).

The executed exact certificate \texttt{code/a\_C4\_fixed.py}
checks (A.37.2), all three harmonic normalizations,
the four branch indices, the cyclic-four jump calculation,
and the homology determinants for all six possible
choices of a two-cusp leaf. Its output is retained
in the adjacent \texttt{.json} and \texttt{.out} files.
\end{proof}
